\exercisegroup

\begin{enumerate}
\def\labelenumi{\arabic{enumi})}
\item
  设 X 是一个林德勒夫空间（TG，第 IX 章，附录 1，第 75 页，定义 1），且 X 是可解圈的。证明，对于每个点 \(x \in X\)，庞加莱群 \(\pi_1(X, x)\) 是可数的。
\item
  设 Y 是一致空间，X 是 Y 的子空间。设 \((\mathrm{U}_{i})_{i\in\mathrm{I}}\) 是 X 的一个开覆盖。
\end{enumerate}

\begin{enumerate}
\def\labelenumi{\alph{enumi})}
\tightlist
\item
  证明，存在 Y 中包含 X 的一个开邻域 V，以及 V 的一个开覆盖 \((V_{i})_{i\in I}\)，满足以下性质：
\end{enumerate}

\begin{enumerate}
\def\labelenumi{(\roman{enumi})}
\item
  对每个 \(i\)，有 \(\mathrm{U}_i = \mathrm{V}_i \cap \mathrm{X}\)；
\item
  对 I 的每个子集 J，若 \(\bigcap_{j\in J}V_j\neq \varnothing\)，则有 \(\bigcap_{j\in J}U_j\neq \varnothing .\)
\end{enumerate}

\begin{enumerate}
\def\labelenumi{\alph{enumi})}
\setcounter{enumi}{1}
\item
  假设 Y 是局部连通的，且各 \(U_{i}\) 连通。证明，存在满足这些条件的覆盖，使得各 \(V_{i}\) 是 Y 的连通开子集。
\item
  假设各 \(U_{i}\) 连通，并对每个 \(i \in I\) 取一点 \(x_{i} \in U_{i}\)。假设 \(\pi_{1}(U_{i}, x_{i})\) 在 \(\pi_{1}(X, x_{i})\) 中的像归结为单位元。证明，存在 Y 中包含 X 的一个邻域 V，使得对每个 \(x \in X\)，典范同态 \(\pi_{1}(X, x) \to \pi_{1}(V, x)\) 存在一个收缩映射。
\end{enumerate}

\begin{enumerate}
\def\labelenumi{\arabic{enumi})}
\setcounter{enumi}{2}
\tightlist
\item
  \begin{enumerate}
  \def\labelenumii{(\arabic{enumii})}
  \setcounter{enumii}{2}
  \tightlist
  \item
    设 \(X\) 是拓扑空间，且 \(x\) 是 \(X\) 的一个点。
  \end{enumerate}
\end{enumerate}

设 G 是一个群，且 \(\varphi\colon\pi_{1}(\mathrm{X},x)\to\mathrm{G}\) 是群同态。设 \((\mathrm{U}_{n})_{n\in\mathbf{N}}\) 是 X 的子空间的递减序列，其交集归结为点 \(x\)。对所有 \(n\)，令 \(\mathrm{G}_{n}\) 为 G 中这样的元素所成的集合：它们是在 \(\varphi\) 下某个包含于 \(\mathrm{U}_{n}\) 中的圈的道路类的像。

\begin{enumerate}
\def\labelenumi{\alph{enumi})}
\item
  若 G 可数，则序列 \((\mathrm{G}_{n})\) 稳定。（归结为对所有 n 都有 \((\mathrm{G}_{n}:\mathrm{G}_{n+1})\geqslant n+1\) 的情形，并取 G 的一个枚举 \((\gamma_{n})_{n}\)；递归构造元素 \(\lambda_{n}\in G_{n}\)，使得记 \(w_{n}=\lambda_{1}\ldots\lambda_{n}\) 时，有 \(\gamma_{i}\notin w_{n}G_{n+1}\)，其中 \(1\leqslant i\leqslant n\)。然后证明交集 \(\bigcap_{i}w_{i}G_{i+1}\) 非空。）
\item
  假设 G 是阿贝尔群，且其单位元是 G 中唯一的无限可除元素。则 \(\bigcap_{n} G_{n}\) 归结为单位元。（设 a 是交集中的一个元素；对每个 n，取 \(\gamma_{n} \in \pi_{1}(U_{n}, x)\)，使其在 \(\varphi\) 下的像等于 a。递归构造道路类 \(\lambda_{n} \in \pi_{1}(U_{n}, x)\)，使得 \(\lambda_{n} = \gamma_{n}\lambda_{n+1}^{d}\)，其中 d 是一个固定的整数且 \(\geqslant 2\)。证明 \(b = (d - 1)\varphi(\lambda_{1}) + a\) 是无限可除的。由此推出 a = 0。）
\item
  更一般地，若 G 是阿贝尔群，且 G 的每个无限可除元素都是挠元素，则 \(G_{n}\) 的交集由挠元素组成。
\end{enumerate}

(3) 本练习及下一练习的结果取自 J. W. CANNON 与 G. R. CONNER 的文章《关于一维空间的基本群》，载于《拓扑学及其应用》153 (2006)，第 2648–2672 页。

\begin{enumerate}
\def\labelenumi{\arabic{enumi})}
\setcounter{enumi}{3}
\tightlist
\item
  设 X 是连通且局部道路连通的拓扑空间。设 x 是 X 的一个点，并且 x 具有可数的邻域基。设 S 是一个集合，且 \(\varphi\colon\pi_{1}(\mathrm{X},x)\to\mathbf{Z}^{(\mathrm{S})}\) 是满射群同态。
\end{enumerate}

\begin{enumerate}
\def\labelenumi{\alph{enumi})}
\item
  证明 S 是可数的。
\item
  假设 X 是紧的；证明 S 是有限的。
\item
  假设 \(\pi_{1}(X,x)\) 是自由阿贝尔群；证明 X 是可解圈的。
\item
  假设 \(\pi_{1}(X,x)\) 是自由群；证明 X 是可解圈的。（利用自由群的下中心列各子群的交集归结为单位元这一事实。）
\end{enumerate}

\begin{enumerate}
\def\labelenumi{\arabic{enumi})}
\setcounter{enumi}{4}
\tightlist
\item
  设 P 是夏威夷耳环（III，第 336 页，练习 6；我们采用其中的记号）。对 \(n \in \mathbf{N}^*\)，令 \(\alpha_n \in \pi_1(\mathrm{P}, o)\) 为 \(\mathrm{C}_n\) 中一个圈的道路类，该道路类生成 \(\pi_1(\mathrm{C}_n, 0)\)。
\end{enumerate}

\begin{enumerate}
\def\labelenumi{\alph{enumi})}
\item
  证明映射 \(\mathrm{Hom}(\pi_{1}(\mathrm{P},0),\mathbf{Z})\to\mathbf{Z}^{\mathbf{N}}\)，\(\varphi\mapsto(\varphi(\alpha_{n}))_{n\in\mathbf{N}}\)，是单射群同态，且其像等于 \(\mathbf{Z}^{(\mathbf{N})}\)。（4）
\item
  群 \(\pi_{1}(P,0)\) 不是自由群；\(\pi_{1}(P,0)\) 对其导出群的商不是自由阿贝尔群。
\item
  设 A 是夏威夷群岛（III，第 338 页，练习 10）。证明从 \(\pi_1(A, o)\) 到 \(\mathbf{Z}\) 的每个同态都是平凡的。
\end{enumerate}

(4) 参见 B. DE SMIT，《夏威夷耳环的基本群不是自由群》，《国际代数与计算杂志》，第 2 卷，第 1 期 (1992)，第 33–37 页。

